\Needspace{8\baselineskip}\originalchapter{Coulomb 碰撞与弛豫}\label{ch:collisions}
\originalsection{碰撞参数与小角累积}
两个带电粒子的相对运动质量为 $\mu=m_1m_2/(m_1+m_2)$, 相对速度 v, 电荷乘积常数 $K=q_1q_2/(4\pi\epsilon_0)$. Rutherford 轨道给 $\tan(\theta/2)=K/(\mu v^2b)$, b 为碰撞参数. 定义 $b_{90}=|K|/(\mu v^2)$, 较大 b 的偏转虽小但靶数很多. 一束测试粒子穿过均匀靶时, b环内碰撞率 $n v 2\pi b\dd b$.

小角横向动量冲量约 $\Delta p_\perp=2K/(vb)$. 因此每单位时间的横速度平方增长为
\begin{equation}
 \frac{\dd\langle v_\perp^2\rangle}{\dd t}=\frac{8\pi nK^2}{m_{
\rm test}^2v}\int\frac{\dd b}b
 =\frac{8\pi nK^2\log\Lambda}{m_{
\rm test}^2v}.
\end{equation}
这里 ``增长'' 的符号应为正. 纵向定向动量减少、横向展宽增加, 二者不应套同一个带负号的定义. 大角散射的非对数贡献在 $\log\Lambda\gg1$ 时是较小修正.

\originalsection{上下截断与 Coulomb 对数}
远处屏蔽或动态响应给 $b_{\max}$, 经典近距离大角偏转及量子衍射给 $b_{\min}\sim\max(b_{90},\hbar/(\mu v))$. 常见静态弱扰动近似取 $b_{\max}=\lambda_D$, 从而 $\log\Lambda=\log(b_{\max}/b_{\min})$. 更快束流的动态屏蔽、磁化碰撞或强耦合可要求不同截断; Coulomb 对数不是不随能量和物种变化的神秘常数.

电子和离子的约化质量与速度分布不同. 同一束流与电子、离子碰撞的 $\log\Lambda$ 可不同; 电子量子截断常比经典 $b_{90}$ 大. 采用 $h$ 或 $\hbar$、速度取 RMS还是 $\sqrt{2T/m}$ 会改变 O(1) 因子, 报数时必须给约定. 在大对数情况下这类差别通常较小, 但不能说所有版本完全相同.

\originalsection{频率、动量损失与热化}
测试粒子在静止靶上的平均纵向动量损失频率, 保留小角对数主项, 为
\begin{equation}\label{eq:fast-slow}
 \nu_{ts}=\frac{4\pi n_sK_{ts}^2\log\Lambda}{m_t\mu_{ts}v^3}.
\end{equation}
对无限重靶, 约化质量接近测试质量, 动量偏转却几乎不改变粒子速率; 因而``纵向束动量消失''不等于``粒子能量都给靶''. 对可动靶, 能量交换需同时跟踪靶的反冲, 多物种热化速率不同.

Maxwell 分布附近小漂移的电子--离子动量频率常用
\begin{equation}
 \nu_{ei}\simeq2.91\times10^{-6}\,
 \frac{n_e[\mathrm{cm}^{-3}]\,Z\log\Lambda}{T_e[\mathrm{eV}]^{3/2}}\ \mathrm{s}^{-1}.
\end{equation}
离子自碰撞频率对应系数约 $4.80\times10^{-8}$、质量因子 $\sqrt{m_i/m_p}$ 和电荷高次幂. 这是指定热频率约定, 不等于每一种测试粒子90度偏转时间. 典型平均自由程 $\ell\sim v_{\rm th}/\nu$ 可远大于装置尺寸, 所以完全电离且稠密不必等于流体局域平衡.

\originalsection{电阻与 runaway}
电子--离子相对漂移承受摩擦, 近似 $\eta=m_e\nu_{ei}/(n_ee^2)$. 对本卷单电荷氢, 修正 Maxwell 电子碰撞后的 Spitzer 平行电阻经验式约为 $\eta\simeq5.2\times10^{-5}\log\Lambda/T_e^{3/2}$ $\Omega\,$m, T用 eV. 一般离子价数还需完整碰撞算子给出的修正系数, 不把本单电荷系数简单乘Z当精确式. 固定对数时密度在分子分母抵消, 但实际上 $\log\Lambda$ 仍缓慢含密度.

快电子拖曳力按 $v^{-2}$ 降低, 外场力 eE近乎不依 v, 所以超过某速度时可能持续加速成为 runaway. 条件比较须使用快粒子拖曳, 不能使用常数热碰撞频率外推到全部能量. 若求得阈值已是几十到数百 keV, 必须检查相对论, 不能一边算到接近 c 一边保持非相对论结论.
